A quantum extreme learning machine (QELM) transforms inputs into quantum
features for a trainable readout without carrying memory between predictions.
Encoding a window of recent inputs allows it to process temporal information.
Can such a window QELM reproduce a quantum reservoir computer (QRC), which
carries information between inputs in a quantum state?
We characterize when this replacement is possible. Two histories with
identical recent inputs are indistinguishable to a window model, yet may
produce different reservoir outputs. For scalar outputs, half the largest
such difference is exactly the smallest achievable worst-case error. Under
suitable regularity assumptions, a growing universal QELM family evaluated
through exact expectations can approach this optimum at every window length.
Arbitrary accuracy is therefore possible precisely when the difference
vanishes uniformly as the window grows. Only
memory that can affect the chosen outputs must fade. The full quantum state
need not forget its past. A two-qubit example demonstrates persistent memory
that prevents arbitrarily accurate finite-window replacement. Simulations
examine the practical roles of window length, feature capacity, and
measurement budgets. Together, these results establish when explicitly
encoded input history can replace recurrent quantum memory and when
information beyond any finite window remains essential.
